\chapter{Introduction}

Collision detection is a concept in Computer Science, which mainly finds application in computer graphics, physics simulations and video games.
Whenever many objects exist in a virtual space, a system must determine which objects intersect or overlap, to simulate interactions between those objects.
A direct solution would be to test each object against each other object in the space, requiring $n(n-1)/2$ pairwise tests for $n$ objects.
For interactive applications, such as video games, where collision detection may occur many times a second, this behavior is often impractical \cite[p.~14]{ericson2004real}.

A common way to reduce the number of required collision checks is to use bounding volumes, which enclose objects with simpler shapes such as bounding boxes.
If two bounding boxes do not overlap, neither will the primitive objects enclosed in them, allowing to skip the otherwise needed exact collision checks.

This principle can be extended through so called Bounding Volume Hierarchies (BVH). 
A BVH organizes objects in a binary-tree, where internal nodes represent bounding volumes containing all their child nodes, and leaf nodes represent primitive objects. 
During traversal entire subtrees can be discarded if their bounding-volumes do not overlap \cite[pp.~235--260]{ericson2004real}.

This work investigates this influence by comparing one representative of each of three general construction mechanisms: top-down, bottom-up, and incremental construction.
The following research question will be addressed by this WAB:

\begin{quote}
\textit{How do Bounding Volume Hierarchy construction strategies influence the number of bounding-volume and primitive collision checks required for all-pairs collision detection in static two-dimensional particle distributions?}
\end{quote}




\newpage